% Shared exercise chunk: l4-copying-context
\begin{exercisechunk}
\item \textbf{Copying with a bounded context.\suppmark}
\label[exercise]{ex:copying-context}
Let \(P\) generate \(B_1B_1B_2B_2\cdots\), where the blocks
\(B_m\) are independent and uniform on \(\{0,1\}^m\).
For an integer \(w\ge0\), let \(\mathcal Q_w\) be all
infinite-sequence models whose next-token distribution can depend
on the position and the last \(w\) tokens only.
\begin{enumerate}[(a),beginpenalty=10000]
\item In a copied block of length \(m>w\), show that the required
bit is fair given the available context, although the full prefix
determines it. Deduce expected conditional KL at least \(\log2\)
at each such position.
\item At \(T=M(M+1)\), for \(M>w\), prove
\begin{equation}
 \inf_{Q\in\mathcal Q_w}\KL(P_{1:T}, Q_{1:T})
 =\frac{\log2}{2}\bigl(T-w(w+1)\bigr).
 \label{eq:copying-context}
\end{equation}
Specify a predictor attaining equality.
\item Explain why more training data cannot remove this error for
a fixed \(w\). How must \(w\) grow to predict every copied bit exactly?
\item Show that an exact recurrent generator must distinguish
\(2^m\) states after generating a fresh length-\(m\) block, and hence
needs at least \(m\) bits of state information.
Compare the description length of the copying rule with its memory
requirement. Relate the distinction to recalling earlier definitions
or intermediate results in a mathematical argument.
\end{enumerate}
\end{exercisechunk}
